\chapter{}
A Sul América do Nordeste, SULANOR, foi inaugurada com uma grande festa para políticos, autoridades, diretores da Sul América Seguros e, Paulo fez questão, para os peões que trabalharam na obra.
Foi um triunfo para Paulo, que vencera o desafio de implantar dois mil e quinhentos hectares de lavoura irrigada e de erguer uma agroindústria de ponta, toda informatizada, naquele lugarejo esquecido à beira do Parnaíba, mas foi principalmente uma façanha daquele povo que, por obra e arte do inusitado empreendimento, foi catapultado do escambo e do extrativismo para uma economia de mercado num passe de mágica.

Os primeiros meses tinham sido desanimadores.
Sob o calor de mais de quarenta graus ao sol, o trabalho não rendia.
Os homens não reagiam sequer ao estímulo de trabalhar em troca de salário e carteira assinada.
Paulo reuniu-os e pediu-lhes uma explicação.
Será que teria que importar trabalhadores de outros estados, quis saber.
Então, um dos peões falou pelos demais:

\textit{``-- Não, doutor, não carece.
É só o senhor fornecer um pedaço de rapadura que seja, para cada um dos homens aqui, porque o trabalho é puxado e a gente só traz de casa água e um punhado de farinha para se sustentar o dia inteiro.''}

Paulo, então, organizou uma cozinha industrial que passou a fornecer diariamente arroz, feijão, carne, verduras e legumes plantados ali mesmo, a todos os trabalhadores da fazenda.
A partir daí, o desempenho daqueles homens cresceu a tal ponto, que nos convencemos ser aquela a melhor mão-de-obra com que jamais havíamos contado.
Rápidos em aprender, versáteis e duros no serviço, viabilizaram no prazo previsto aquela obra que para muitos não passava de um delírio conjunto dos donos da Sul-América e de um governador visionário.
Para dar uma ideia da dimensão do que ali aconteceu, é suficiente dizer que o peão que controlava com perícia os sofisticados comandos eletrônicos da usina de beneficiamento e empacotamento do arroz, era analfabeto.

\begin{figure}
\centering
\includegraphics[width=0.9\linewidth]{37/implantação-sulanor.png}
\caption{A implantação da SULANOR.}
\end{figure}

\begin{figure}
\centering
\includegraphics[width=0.9\linewidth]{37/supervis.png}
\caption{Supervisionando o plantio de arroz.}
\end{figure}

\begin{figure}
\centering
\includegraphics[width=0.9\linewidth]{37/brincando-usina.png}
\caption{Brincando na usina de arroz.}
\end{figure}

A SULANOR bateu recordes de produção, mas, ao contrário do que se podia esperar, começou a apresentar sucessivos prejuízos.
Misteriosamente, as vendas do arroz realizadas pela equipe do Paulo, com excelentes margens de lucro, eram anuladas pelo escritório central do Rio de Janeiro que preferia entregar o produto a preços aviltados.
A prática se repetiu com tamanha regularidade que pareceu claro a todos que a coisa era intencional.
A fazenda devia dar prejuízo.
Tratava-se, de novo, daquela prática de usar as empresas rurais constituídas a partir de estímulos fiscais para mascarar os habituais desvios de caixa das respectivas matrizes.
A partir dessa desconfiança, as relações entre Paulo e seus superiores começaram a azedar.
A gota d'água foi uma reunião, no Rio de Janeiro, em que a empresa sugeriu uma reforma no contrato de trabalho do Paulo que, na prática, implicaria em redução de salário.
Paulo demitiu-se.

Meu trabalho no Diocesano, além do prazer, nos permitia manter os cinco num colégio particular graças às bolsas de estudo e descontos a que eu tinha direito como funcionária da casa.
Tinha mais facilidade em obter livros a preços menores ou até de graça, às vezes.
Mas, o salário apenas cobria o pagamento das moças que me ajudavam em casa, para que eu pudesse trabalhar.
Daí que Paulo teve que procurar o que fazer.

De Conceição trouxéramos o Sr. Pedro, um japonês que cuidava do viveiro de mudas de seringueira e tinha se revelado um formidável horticultor.
Numa chácara arrendada, Paulo montou uma horta para fornecer verduras e legumes na cidade.
O viveiro de seringueira também foi reorganizado.
Paulo até conseguiu um cliente, um empresário pernambucano que se entusiasmou por produzir borracha na sua fazenda de Amarante, no interior do Piauí.
Até hoje, tantos anos passados, o bosque de seringueiras de Amarante desperta curiosidade e admiração.
Ou melhor, que restou dele, porque o empresário morreu poucos anos depois, a fazenda foi abandonada e um punhado de sem-terras acabou por invadi-la, pondo fogo na maior parte das árvores, ninguém sabe por quê.

Vender verduras e legumes em Teresina daquela época não era fácil, embora fornecêssemos os maiores supermercados da cidade.
É que os piauienses não tinham o hábito de consumi-los em quantidade.
Principalmente verduras.
E os legumes eram comprados por unidade.
Por mais de uma vez, antes de plantarmos a horta, perguntaram-me no mercado onde eu fazia a feira semanal, se eu tinha pensão ou hotel, porque eu pedia quilo disso, dúzia daquilo, maço daquele outro.
Era um espanto! E a horta do seu Pedro produzia às toneladas.
Houve um dia que a nossa garagem ficou cheia de alto a baixo de maravilhosos pepinos que não tínhamos para quem vender.
E o supermercado dizia que cada berinjela e pimentão daqueles excedia o peso que os fregueses estavam dispostos a levar, pois não davam conta de consumi-los antes que se estragassem.

Por essa época, chegou ao conhecimento do Paulo que uma velha senhora da sociedade teresinense queria arrendar suas terras à beira do Parnaíba.
Paulo foi procurá-la e o negócio ficou acertado.
A terra era toda mecanizável, mas precisava de recuperação.
Paulo decidiu fazer um cultivo de melancias irrigado, pois os piauienses eram loucos pela fruta e ela tinha que vir de muito longe.
Depois plantaria soja, se me lembro bem.

Certa manhã, ele e Carlito, o motorista que o acompanhava desde os tempos da SULANOR, saíram muito cedo, carregando sobre a nossa caminhonete, à qual foi engatada uma carreta, os cerca de seiscentos canos de irrigação que compunham nosso equipamento.
Eu saí em seguida, para trabalhar.
Por volta das onze horas, recebi um inusitado telefonema do meu compadre Chico que, com voz tensa, pediu que fosse ao encontro dele no Hospital Getúlio Vargas, onde ficava o maior pronto-socorro público de Teresina.
Paulo tinha sofrido um acidente na estrada.


Quando, chegando à emergência, vi Paulo amontoado sobre a maca, senti o mundo desabar.
O acidente tinha sido muito grave.
Um pneu traseiro da caminhonete estourara e Carlito perdera a direção, descendo um declive de sete ou oito metros ao fim do qual veículo e carreta se arrebentaram contra uma palmeira de babaçu.
Sentado no banco do carona, Paulo recebera todo o impacto da batida, sofrendo fraturas cruzadas do ombro direito e da cabeça do fêmur esquerdo.
E uma pancada na cabeça que o poria em estado de confusão por meses.
Ambas as articulações estilhaçaram.
Sem próteses disponíveis para operá-lo, dois heroicos ortopedistas manipularam os cacos de ossos na tentativa de recompor razoavelmente o ombro e o quadril.
Depois, mandaram-no para o quarto e o alçaram a um aparelho que mais parecia um instrumento de tortura medieval, dotado de pesos que estiravam a perna que se retraíra abdômen adentro.
Ali ele ficou por um mês.

Afastei-me do Colégio para cuidar do Paulo.
João, já médico residente, veio a Teresina para me dar força e avaliar o que tinha acontecido.
Não sei o que contou ao voltar, mas os telefonemas que eu recebi, da minha família e da família do Paulo, revelavam muita preocupação.

O mais difícil de tudo, é que eu nunca tivera muito tempo de me pôr a par de detalhes dos negócios do Paulo.
Aliás, nunca foi muito do estilo dele dar satisfações minuciosas do que faz.
E o estado de confusão em que ele ficou tornava quase impossível entender o que eu devia ou não devia fazer.
Eu pegava o carro e corria à chácara e à fazenda, para providenciar o que os empregados pediam, fazia pagamentos, mas não podia me afastar de casa por muito tempo e nem avaliar com cuidado a necessidade ou não daqueles gastos todos.
O fato é que ao longo dos seis meses que se seguiram ao acidente, tudo o que tínhamos conseguido juntar se foi, pouco a pouco.

Paulo veio para casa engessado das axilas aos pés.
As crianças e eu tivemos que improvisar uma espécie de guindaste nas grades da janela do quarto para levantá-lo da cama, como um grande tronco, para poder fazer-lhe a higiene e trocar os lençóis.
Era um mutirão que eles encararam com o habitual entusiasmo.
Da cama, um mês e meio depois, ele foi para a cadeira de rodas.
Aí eu já podia conduzi-lo para a fazenda e para outros lugares e ele começou a tentar assumir algumas decisões, mas eu percebia-lhe a mente enevoada.
Vieram as sessões de fisioterapia, mais dois longos meses e, por fim as muletas.

Numa manhã, inesperadamente meu diretor assomou à porta do quarto onde eu cuidava do Paulo.
Viera pedir que eu reassumisse minhas funções porque minha substituta andara metendo os pés pelas mãos e a Coordenação Pedagógica estava um caos.
Paulo me liberou e, a partir daí, foi à luta sozinho, de muletas, mancando ainda e, tenho certeza, sentindo um bocado de dor.

O acidente do Paulo marcou o início de uma fase dificílima da nossa vida e do nosso casamento.
Era bem verdade que no que diz respeito a negócios, ele e eu nunca estivemos muito de acordo, de modo que esse era um assunto espinhoso que eu sempre preferi evitar.
Mas, a partir do acidente, dadas as dificuldades que se seguiram, ficava cada vez mais complicado fingir que não me dava conta do que acontecia.

